% frontmatter/notation.tex
\chapter{符号与读数约定}

本表给出全书主线的符号。数学符号是帮助读者追踪对象的约定，不是对象已经被准确测量的证明。同一字母在不同论文中可能具有不同含义；引用外部公式时，本书保留其局部定义，并在发生冲突时说明。尤其不要把一个版本编号、一次任务尝试和一次模型请求互相替换。

\begin{longtable}{@{}p{2.4cm}p{10.5cm}@{}}
\toprule
符号 & 含义 \\
\midrule
\endhead
$t$ & 任务内的步（step） \\
$e$ & 一次任务尝试（episode）；其轨迹另记为$\tau_{k,e}$ \\
$k$ & 主线的系统版本索引；技能成本例题中也可表示复用次数，以该节定义为准。 \\
$\theta$ & 模型参数 \\
$P$ & 指令与持久策略 \\
$M$ & 经验记忆 \\
$L$ & 工具与技能库 \\
$G$ & 工作流与角色结构 \\
$\mathcal{V}$ & 验证器（验收器、任务划分、最终测试集）；\textbf{默认冻结} \\
$A_k$ & 系统描述$(\theta_k,P_k,M_k,L_k,G_k;\mathcal V)$：五类可更新组件与单列的评价设置。 \\
$B$ & 总预算（token、调用次数或金额，须写明单位） \\
$Q$ & 单个候选版本查询选择验证集的题次上限 \\
$r$ & Gate 的灰度放量比例 \\
$\Delta_{\text{migrate}}$ & 换模型或换 harness 后的迁移收益 \\
\bottomrule
\end{longtable}

\section*{三个时间尺度不要混记}

任务内步骤描述一次求解怎样推进，例如读取文件、生成候选和运行公开测试；任务尝试把这些步骤组织成一个可以验收的提交；版本层描述会被后续任务继承的状态如何变化。本书将进入后一层的持久变化作为自进化研究对象。这是本书的操作性约定，不要求所有外部论文也使用完全相同的术语。

一次任务尝试并不必然只调用模型一次。一个重试策略还可能在同一道题上产生多个候选，最后按公开信息选一个提交；报告这样的成功率时，要同时给候选数量、选择规则和资源，不把它写成基础模型的一次生成成绩。类似地，一次版本更新可能包含生成、筛选、回归和发布多个动作，不能只按最后一次写文件的成本计算。

评价设置默认冻结是比较规则，不是“验收程序永远不能修错”的要求。若发现测试错误，可以保留旧版本和受影响记录后修正，并为修正后的比较建立清楚的版本边界。旧评价与新评价的分数不能无说明地拼接。若研究目标本来就是评价器更新，还需把它设为独立干预，并另行回答新评价由谁检验。

\section*{差值、百分比与符号方向}

若旧成功率为\calcnum{0.50}、新成功率为\calcnum{0.60}，两者相减为\calcnum{0.10}，也就是提高十个百分点；以旧成功率为分母的相对提升是\calcnum{20\%}。这是教学算术约定，不能把“十个百分点”和“百分之十”互换。写“提升两倍”也应说明究竟指新值为旧值的两倍，还是增量为旧值的两倍。

在本书的退化分析中，$D=\text{旧版本得分}-\text{新版本得分}$，所以$D>0$表示下降；一般收益差则常按新减旧定义，正号表示改善。符号方向必须随表格、图注和代码字段一起核对。区间不含零但全部为负，不能支持预先写成正向的命题；区间跨零也不能直接改写成两个方案等效。

预算同样需要附上单位。相同调用次数上限并不等于相同实际token、墙钟或费用；未知计费用量不能填零。当前版本将已发生的实际总账与假设独立部署某方案的成本分开，避免共享准备费被重复累加，也避免把共享材料当成免费资源。读成本图之前，应先找到它使用的是哪一本账和哪个任务序列长度。

\section*{公式中的数据身份}

正文把外部来源报告的读数、本书真实运行的读数和根据指定输入得到的计算分别标为引用、实测和计算。一个数据表可以同时包含这三类内容，但每一项应能回到相应来源。闭式公式中的参数、教学设定中的小数、规划样本量都不能因为出现在实验章就变为实测。

理论推导还应附有条件。例如，在独立误差下得到的方差加法，不会自动覆盖同一任务流中共享记忆造成的依赖；给定优势估计下的最优策略，也不会自动保证真实任务表现改善。阅读时先确认公式描述的是定义、条件命题、近似还是经验观察，再看其数值。这样既保留公式的作用，也避免把符号的精确外观误当成证据的完整程度。

第16章将目标版本的评分规则另记为$\mathcal O_0$、$\mathcal O_1$，以免与工作流结构$G$混淆；比较两版系统时，应说明沿用哪一版规则。

另外，$t$在任务流分析中可能重新定义为任务到达索引，$r$在统计例题中可能表示重复索引。局部定义不把全局放量比例变成随机变量，也不把版本索引变成任务内步数。公式连接到代码之前，应写出它所用的全部索引范围和数组维度；遇到字母相同但对象不同的两式，先改名对齐再推导。
